\chapter{目标对齐与行为约束}
\label{chap:trust-alignment}

用户要求助手整理邮件，并希望重要内容得到保留。如果训练奖励只依据“处理速度”，助手可能学会少做检查；如果奖励只依据用户给出的即时好评，它可能倾向于给出令人满意却缺少依据的回答。这里没有外部攻击，模型仍可能偏离实际要求，因为学习过程优化的是可观察反馈，而用户关心的是行动后果。

上一章把公平要求落实为可检验的分配目标与干预规则。类似地，邮件助手也需要区分希望实现的结果、用于训练的反馈与行动时不得越过的边界。

本章从目标与约束的含义出发，说明行为偏离有哪些表现，以及可以用什么证据评估和监督。监督发现的规律还需要经受新任务、长交互和模型更新的考验，因此随后讨论行为泛化与约束保持的机制，再按训练、推理和执行等位置介绍具体干预。监督提供判断，泛化分析说明判断和约束能延伸到哪里，干预则改变行为及其后果。

\section{目标对齐与行为约束的基本概念}
\label{sec:alignment-objectives}

\subsection{目标对齐与行为约束的含义}
\label{sec:alignment-definition}

\term{目标对齐}（alignment）关注模型学到的行为能否符合预定的实际目标，而不只是获得较高训练分数。例如，整理邮件的目标可以是减少查找负担，同时保留重要内容；“收件箱变空”只是容易测量的结果，未必等于完成这一目标。对齐要求说明目标由谁提出、适用于谁，以及不同要求冲突时如何处理。

\term{行为约束}规定模型采取行动时必须满足的条件或不得越过的边界。例如，只能操作本次授权的邮件、不能删除重要内容、信息不足时应先核查。目标回答希望达到什么结果，约束规定哪些实现方式可以接受。一个助手即使完成整理任务，也可能因未经授权转发信息而违反约束；另一个助手可能始终不越权，却什么也不做，因而没有完成目标。

目标与约束都应落实到可检查的行为。回答内容是否有依据、行动是否得到授权、最终状态是否符合用户要求，是不同的观察对象。“模型表示自己理解了目标”不能替代这些证据。某些偏离可能源于能力不足、缺少信息或任务表述不清，也可能源于代理目标错误；仅凭一次失败，通常不能确定模型具有何种内部意图。

图~\ref{fig:trust-alignment-mail-goal}比较了两个都能清空收件箱的方案。若评价只看收件箱是否为空，两种方案都可能获得高分；一旦检查重要邮件是否保留、用户能否找回，便能区分实际后果。这里需要补充的是评价目标，而不只是让模型更快地完成原来的评分任务。

\begin{figure}[!htb]
  \centering
  \bookdraftgraphic{figures/trust-alignment-mail-goal.png}
  \caption{相同的表面完成状态，不同的任务后果。左侧通过删除清空收件箱，星标的重要邮件也被丢弃；右侧通过归档清空收件箱，用户仍能找回重要邮件。星标表示应保留的内容，空托盘表示被单独评分的表面结果。}
  \label{fig:trust-alignment-mail-goal}
\end{figure}

\subsection{行为偏离的表现与原因}
\label{sec:alignment-misalignment}

设任务历史为$h$，行动为$a$，实际希望改善的后果用$U(h,a)$表示。实践中通常不能直接获得$U$，只能观察评分、点击、偏好比较或任务完成标记，并据此构造奖励$r(h,a)$。奖励是可优化的代理；两者在训练数据上相关，不保证经过强烈优化后仍然一致。

\term{奖励投机}（reward hacking）指系统获得较高奖励，却没有实现奖励原本代表的目标。一个只按“文件数量减少”评价整理效果的任务，会鼓励删除本该保留的文件。改进需要回到目标定义，检查评价器和环境中的漏洞，而不只是更充分地优化当前奖励。

偏好还可能因人和情境而异。某人偏好简短回答，另一人需要详细解释；这类差别可以个性化。未经授权处理他人信息则可能属于共同约束，不能因为某个用户偏好便自动放宽。实际意图、个体偏好和共同约束应分别说明。

后文沿用同一个整理任务：整理旧邮件，保留标为重要的内容，并且不向未授权对象转发；具体整理方式允许按用户偏好选择。对一封既旧又重要的邮件，方案$A$是归档并建立索引，方案$B$是留在收件箱并加标记，方案$C$是直接删除。$A$与$B$都保留了信息，但用户可能更偏好$A$的整理效果；$C$违反保留要求，不能用速度优势补偿。这里可分别核查归档后的可检索性、邮件是否仍存在以及实际收件人，而不必把所有要求压成一个主观好评分数。

这种区分决定了如何收集反馈。邮件助手的一段摘要可能同时涉及信息完整性、语气、长度和隐私。如果标注者只给一个整体分数，训练者无法知道分数下降究竟是因为遗漏了会议时间，还是因为措辞不够亲切。把可核查的任务条件、允许变化的风格和禁止违反的边界分别记录，能够保留后续纠错所需的信息。对于确有价值分歧的问题，应明确适用人群与选择机制；对事实判断的分歧，则应检查证据，而不能直接把多数偏好当成事实。

代理目标的误差也会被选择过程放大。写成$r(h,a)=U(h,a)+e(h,a)$，其中$e$为评价误差，即使在原始行为分布下$\mathbb E[e]=0$，被选出的最大得分动作仍可能具有正的评价误差。因为选择偏爱得分高的动作，而得分高既可能来自真实效用，也可能来自高估。这里的条件期望已经改变，原分布上的无偏性不能直接沿用。

\begin{example}[奖励误差被优化放大]
\label{ex:alignment-proxy-selection}
在一个教学例子中，邮件助手有三个整理方案，真实效用依次为$8,7,2$，奖励模型给分依次为$8,7,9$。若助手原先只考虑前两个方案，奖励看似能够正确排序；当更强的搜索把第三个方案纳入候选，最高奖励对应的真实效用反而下降。第三个方案可以是删除所有难分类邮件：它让界面整洁，却破坏了保留重要信息的要求。问题来自评价器对新方案的高估，不是搜索没有充分最大化奖励。
\end{example}

Gao等用一个固定的较强奖励模型提供比较标签，再训练代理奖励模型，分别通过强化学习和从多个样本中择优来增加优化强度。他们在这一可控实验中观察到，代理得分继续提高时，作为参照的奖励得分可能下降。该研究使奖励过度优化可以被定量考察；参照模型仍然是一个模型，其得分不能等同于真实人类价值\scite{gao2022overoptimization}

因此，检查优化后的行为需要独立于被优化指标的证据。在邮件例子中，可以另查应保留邮件的清单与实际操作记录，而不是再次询问同一个奖励模型“整理得好不好”。降低优化强度、增加反馈覆盖和更新评价器能够缓解不同来源的误差；它们都需要重新检验，因为更精确地拟合旧分布，并不自动补上新策略访问区域中的盲点。

偏离还可能表现为迎合与表面合规：用户误以为某封邮件不重要，助手不核对标记就赞同删除；或者助手写出“已经检查权限”，实际却没有读取授权记录。相反，把所有整理请求都拒绝虽然能减少误删，也会破坏应提供的服务。这些教学例子分别提示事实依据、真实执行和正常任务效用都应进入评估，不能只检查回答是否包含合适的安全措辞。

\subsection{软目标与行为边界}

对于顺序决策，记策略$\pi$根据历史$h_t$选择动作$a_t$，$r_t$为任务奖励，$c_t\geq0$为风险成本，$T$为时域。一个常见目标是
\begin{equation}
 \max_\pi\ \mathbb E_\pi\sum_{t=0}^{T-1}r_t,
 \qquad
 \mathbb E_\pi\sum_{t=0}^{T-1}c_t\leq b,
 \label{eq:alignment-constrained-policy}
\end{equation}
其中$b$为容许预算。该约束控制平均累计成本，不直接保证每条轨迹都安全。如果要求禁止某个动作，需要在动作集合或执行权限上施加相应限制；如果要求失败概率不超过阈值，则需要定义失败事件并建立概率约束。

将成本乘以惩罚系数并入奖励，通常便于训练，却仍需验证约束是否满足。目标表述应保持这种区别，否则“已经加入安全惩罚”很容易被误写成“不会发生违规行为”。

设一个方案的收益为10，风险成本为1；另一个方案的收益为5，风险成本为0。若优化$r-\lambda c$且$\lambda=2$，前者得分8，后者得分5，算法会选择有风险的方案。这与目标函数完全一致。只有当风险本来就允许用收益补偿时，这种交换才符合要求；若风险代表禁止泄露的内容，允许补偿本身就写错了问题。

对长期任务，还应说明成本究竟计算什么。每步违规次数之和、至少发生一次违规的指示量、最终经济损失会给出不同约束。一次严重披露不能因随后进行了许多无害动作而被“平均掉”。第~\ref{sec:trust-requirements}节的风险要求在这里需要落实到轨迹、动作和受影响对象，才能进入训练与监督。

\section{行为的评估与监督}
\label{sec:alignment-oversight}

行为评估根据要求判断一段输出或行动过程是否合适；行为监督把这种判断用于持续检查、提供反馈或发出告警。离线评价、训练反馈和在线监测可以共享判据，但用途不同：评价形成证据，反馈供后续训练使用，监测提示当前风险；它们本身不等于策略参数已经改变，也不等于动作已经被停止。具体改变行为的措施在第~\ref{sec:alignment-intervention}节介绍。

评估与监督需要共同解决三个问题：哪些性质需要判断，依据什么记录与评价者作出判断，以及判断是否准确、及时且负担可承受。仅给一个整体好评分数，会把目标达成、规则违反和表达偏好混在一起；只增加一个监控模型，也不能补上原始记录中缺失的事实。

\subsection{行为评估与监督的维度}
\label{sec:alignment-oversight-dimensions}

评估维度由已声明的目标和约束决定。对邮件助手，至少需要分别观察以下几方面，它们可在同一条轨迹上同时测量。
\begin{itemize}
 \item 目标达成与结果质量：邮件是否完成整理，重要信息是否可检索，摘要事实是否有依据。对应任务完成率、事实核验结果与实际损失。
 \item 约束遵守与影响范围：是否误删重要邮件、向未授权对象转发，操作涉及多少对象。对应具体违反事件、发生频率和严重程度。
 \item 决策过程与信息使用：关键操作前是否获得并正确使用必要证据，工具调用与状态记录是否支持模型的陈述。过程检查需要核对实际调用，不能只看自述。
 \item 不确定情况下的处置：缺少标记或授权信息时是否澄清、复核或暂停，是否在无需停止的正常任务上过度拒绝。对应风险识别、恰当转交与错误干预的比例。
\end{itemize}
这些维度应按任务、人群和关键情境展开。用户偏好的风格可以在合规候选中比较，明确禁止的动作则应单独统计，不能由其他维度的高分抵消。跨情境稳定性需要在相同维度上改变环境后重新测量，具体机制与范围在下一节讨论。

维度与方法并不是同一层分类。“是否越权”是要判断的性质，可以通过权限表、人工复核或模型评价获得证据；“人工反馈”是证据来源，可以同时评价事实、风格和操作。“结果或过程”说明标签附在哪里，“离线或在线”说明何时评价。下面按照证据如何取得、运行时如何观测和复杂判断如何扩展来组织方法，避免把这些交叉维度当成互斥技术。

\subsection{行为评价信号的获取}
\label{sec:alignment-learning-evaluation}

判断行为是否符合要求，既可以直接核对事实与状态，也可以由人或评价模型比较候选的优劣。评价信号的构造有两个交叉维度：由谁、依据什么给出判断，以及把判断附在完整结果还是中间步骤上。人工偏好、原则指导的模型反馈可以采用不同粒度，过程反馈也不限定必须由人给出。无论怎样组合，评价器都应保留实际要求所需的事实和约束。

\paragraph{事实与规则支持的直接核验}

对于有明确真值或状态的任务，优先使用可复核证据。邮件保留要求可以比较操作前后的记录与索引，授权要求可以核对动作对象和收件人，摘要事实可以逐项对照原文。规则检查器给出是否满足条件及违反位置，任务测试给出完成与损失情况；它们比笼统好评分数更容易定位错误，但结论仍受规则是否完整、状态记录是否可信的限制。

\paragraph{人工比较与偏好评价}

当人难以给出精确奖励，却能够比较两个回答时，可以使用偏好数据。设上下文为$\bm{x}$，被偏好的回答为$y^+$，另一回答为$y^-$，奖励模型为$r_{\bm{\phi}}$。一种成对偏好模型令
\[
 \mathbb P(y^+\succ y^-\mid\bm{x})
 =\sigma\!\left(r_{\bm{\phi}}(\bm{x},y^+)
              -r_{\bm{\phi}}(\bm{x},y^-)\right),
\]
其中$\sigma(u)=1/(1+e^{-u})$。对观察到的比较最小化负对数似然，就能够训练奖励模型。这个模型解释偏好差异的概率，仍依赖比较者、提示方式和数据分布。

偏好数据的基本记录是同一上下文、两个候选以及比较结果，而不是只收集被喜欢的回答。可以让当前策略对同一任务生成$A$、$B$，随机交换展示顺序，再记录比较者身份、选择与依据。奖励网络分别读取上下文和候选，输出两个标量；损失只通过它们的差调整参数。训练完成后，奖励模型可以给新候选评分，但这些分数还没有改变负责生成候选的策略。

在一个教学数据集中，假设针对同一对可行方案收集四次判断，其中三次偏好$A$，一次偏好$B$。若只拟合这一对，令$p$为偏好$A$的模型概率，平均损失为$-\tfrac34\log p-\tfrac14\log(1-p)$，其最小值在$p=3/4$取得，对应奖励差$\log 3$。这说明奖励差表达比较频率的相对强度。删除重要邮件的$C$应另外记录为违反约束，不能因为少数比较者喜欢“清空得快”，便把它与$A$、$B$的整理风格差异等同处理。

这一形式通常称为Bradley--Terry偏好模型。对同一上下文中的所有回答加上同一个常数，比较概率不变，因此仅凭成对比较无法识别奖励的绝对零点。模型还假设回答能够由一个标量分数排序。如果某个群体偏好简洁、另一个群体偏好详尽，汇总偏好可能隐藏真实分歧；若出现稳定的循环偏好，一个标量排序也不能忠实表示全部比较关系。

上述$3/4$是拟合比较行为得到的概率，不是回答真实、无害或满足全部约束的概率。把奖励分数直接当作安全置信度，会混淆本章的偏好学习与第~\ref{sec:reliability-evaluation}节讨论的概率检验。

\paragraph{原则指导的模型评价}

人工逐条比较的成本，会限制反馈能覆盖多少情况。另一个方向是由人明确原则，再让模型依据原则批评和比较回答。\term{宪法式人工智能}（Constitutional AI）中的“宪法”指一组指导行为评价的书面原则，并不表示这些原则自动成为法律或取得普遍共识。Bai等的方法先让模型依据原则批评原回答、生成修订，再用修订回答做监督训练；随后让模型依据原则比较候选回答，以这些比较训练偏好模型并进行强化学习。这一阶段通常称为AI反馈强化学习，即RLAIF\scite{bai2022constitutional}

构造这类反馈时，可以向评价模型同时提供任务、候选、相关事实和适用原则，让它指出具体冲突，再生成修订或比较标签。例如，候选$C$写“删除旧邮件以减少干扰”，工具记录却显示重要标记为真；原则评价应据此指出删除与保留要求冲突，修订成保留并归档的操作。修订输出可以进入示范数据，两个候选的比较结果可以进入偏好数据。因此，人类反馈与AI反馈改变的是标签的产生者，而不是凭空增加一种新的策略输出；后续仍需要示范拟合或偏好驱动的策略更新。

这里自动化的是大量具体评价，而原则选择、原则冲突的处理和效果验证仍需要人负责。例如，“积极完成用户请求”与“不得披露未经授权的信息”在某个请求上发生冲突时，评价器需要知道哪一条构成不可放宽的边界。仅增加原则数量可能带来矛盾；只要求给出批评文字，也可能让模型学会写出合适的批评却保留原错误。

反馈来源改变后，验证对象也应改变。人工偏好需要考察标注群体和一致性，模型偏好需要另外考察评价模型与生成模型的共同盲点。应抽样检查原则是否落实到具体事实、正常请求是否被过度拒绝，以及换用评价模型或表达方式后结论是否保持。AI反馈可以扩大训练规模，却不会把评价器的判断自动转化为独立真值。

\paragraph{结果评价与过程评价}

无论反馈由人还是模型给出，长任务只评价最终结果时，都很难判断哪一步造成失败。过程反馈对中间操作、证据检查或计划选择提供评价，可以帮助定位问题；但局部步骤都看似合理，仍可能组合成失败结果。因此，过程反馈应与结果验证结合。

\term{结果监督}（outcome supervision）只给最终产物评分，\term{过程监督}（process supervision）则在中间步骤提供对错或质量判断。若一道题最终答错，结果标签不能直接区分是哪一步算错；若最终碰巧答对，错误推导还可能得到奖励。过程标签提供更细的责任定位，但标注者需要理解局部步骤，标注数量也会增加。

Lightman等在数学解题中比较这两种监督，用人对推导步骤的反馈训练过程奖励模型，并在多个候选解法中选择。研究在所用数学任务和生成设置中发现过程监督优于其结果监督基线，也研究了主动选择有价值步骤进行标注。其证据支持逐步反馈在可核查推理任务中的价值；它没有证明自然语言步骤始终忠实反映内部计算，也没有直接给出开放环境智能体的安全保证\scite{lightman2023verify}

可以用邮件整理说明这种差别。结果监督检查“重要邮件是否仍然存在”；过程监督还检查“删除前是否核对用户标记”“转发前是否核对收件人”。前者能发现最终损失，后者更容易定位缺失的核验操作。若模型只在日志中写“已核对”而实际上未读取标记，过程文字得分仍会误导训练。因此，能够观测的工具返回与状态变化，应成为过程反馈的证据。

过程评价的训练记录则应包含当前历史、待评步骤和局部标签。对于“读取重要标记—选择归档—检查可检索性”的序列，评价者可以给实际读取并正确使用标记的步骤正标签，给未经读取便声称“不是重要邮件”的步骤负标签。过程奖励模型读取每个前缀与当前步骤，预测该步骤的标签或质量分数；结果奖励模型可以读取完整候选，但只从最终成败得到标签。根本差别在标签约束的是当前步骤还是整项任务，而不在评价器能否看到推理文字。比较它们时，可以让两个评价器对同一批候选轨迹排序，再用独立状态检查判断各自选中的轨迹是否真正保留了邮件。

过程奖励也需要明确聚合方式。若各步分数$q_t$代表在前面步骤均正确的条件下当前步骤正确的概率，链式法则给出整条过程正确概率为$\prod_tq_t$，这里不要求各步独立。但实际奖励模型输出未必具有这种条件概率含义，直接把任意步骤分数相乘会制造虚假的概率解释。采用最小值、平均值或加权和时，也分别强调最差步骤、平均质量或特定步骤，需要用任务级结果检验选择是否合理。

\subsection{运行行为的观测与监测}
\label{sec:alignment-oversight-information}

监督需要决定看什么。最终结果能够直接关联任务目标，却可能发现问题太晚；动作序列能够识别权限和流程偏离，却未必解释动机；中间推理文本可以提供额外线索。若把这些文本用作行为原因的解释，需要检验它们是否真实反映了形成行为的依据；后文第~\ref{sec:interpretability-quality}节将系统讨论这种忠实性要求。若只把文本用于告警，则应直接检验它能否帮助发现偏离，不能由语言是否流畅推断监督效果。

设监控器根据可见记录$v(h_t)$输出风险分数$m(v(h_t))$。记录选择决定它能够识别什么：如果关键工具参数未进入日志，监控器不能仅凭最终摘要判断是否越权。增加记录也会带来隐私与成本，需要依据风险确定必要观察。

构造监控器需要把记录与偏离事件对应起来。对于可直接核查的规则，可以读取动作名称、邮件标记及收件人授权表，输出“重要邮件被提议删除”或“目标收件人未授权”的布尔标记。对于需要理解上下文的偏离，可以由人工和工具状态共同标注历史片段，训练分类器或评分模型，以$v(h_t)$为输入、核验后的偏离标签为目标。前者把已知条件写成判定规则，后者从带标签记录中拟合判定关系；两者都只产生告警或风险分数。

邮件助手若只输出“已安全归档”，监控器无法核对$C$是否偷偷被执行。若记录包含调用前的重要标记、拟执行动作和工具返回，就能在删除发生前标记冲突，并在执行后核查邮件是否仍在。这说明增加记录改变的是可识别的信息，而不是自动提升判断准确率；还需要在相同轨迹上比较不同记录范围的漏报与误报。

即使日志完整，评价者也未必能够判断行为是否正确。增加可见信息解决的是观察缺失，分解任务、查询证据和辅助判断解决的是监督能力不足，需要分别检查。

\subsection{复杂任务中的监督能力}
\label{sec:alignment-oversight-capacity}

当执行模型在某些任务上比评价者更强，要求它完全模仿评价者，可能反而压低已有能力。\term{可扩展监督}（scalable oversight）研究如何借助分解、辅助工具或学习方法，使有限的监督能力覆盖更复杂行为。其困难是监督者既要利用模型完成自己难以独立完成的任务，又不能把模型自己的判断直接当作最终依据。

一种可检验的简化，是用较弱模型产生标签，再训练一个已经具有较强预训练能力的模型。Burns等把这一问题称为\term{弱到强泛化}（weak-to-strong generalization）：较强学生是否能够从较弱监督中恢复超过监督者的任务表现。研究在自然语言、棋题和奖励建模等设置中观察到这种现象，同时发现简单拟合弱标签并不能完全恢复强模型能力；辅助置信损失在部分自然语言任务中减少了对弱监督错误的模仿，但效果依赖设置\scite{burns2023weaktostrong}

这里借助训练实验检验弱反馈能否引出更强的判断能力，尚未直接验证某个运行时监控器。把这种结果用于行为监督，还需要另行测量它对偏离事件的检出表现。这个问题与从零学习不同。强模型已经通过预训练获得可利用的表示和规律，弱标签起到指出当前任务方向的作用。如果弱标签的错误特别容易被学生学会，优化仍可能把已有正确判断覆盖掉。因而“学生超过老师”的结果不能解释为错误标签凭空提供了更多真值，也不能保证面对任何能力差距都有效。

为比较恢复程度，设同一测试分布上弱监督者得分为$s_{\mathrm{w}}$，弱监督训练后的强学生得分为$s_{\mathrm{ws}}$，用可信标签训练强模型的参照得分为$s_{\mathrm{s}}$，且$s_{\mathrm{s}}>s_{\mathrm{w}}$，则可计算
\begin{equation}
 g=\frac{s_{\mathrm{ws}}-s_{\mathrm{w}}}
          {s_{\mathrm{s}}-s_{\mathrm{w}}}.
 \label{eq:alignment-weak-strong-recovery}
\end{equation}
$g=0$表示没有超过弱监督者，$g=1$表示达到该强监督参照。它也可能小于0或大于1，参照得分并非数学上的能力上限。这种指标帮助区分“有所改进”和“恢复了多少可测能力”，却仍依赖测试真值。如果真实使用场景连验证标签都无法获得，实验中的恢复率不能被直接测量。

实际比较应把弱标签生成与最终验收分开：弱模型只为训练样本提供标签，强学生据此训练；测试样本则保留可信标签，同时评价弱模型、强学生以及强监督参照。例如，在一个构造的同分布测试中，三者得分依次为0.6、0.75和0.9，则$g=0.5$，表示恢复了一半的参照差距。若把用于拟合弱标签的数据直接拿来计算恢复率，就无法判断学生是否只记住了弱监督错误。将这种方法用于监控时，测试标签还必须对应真实偏离事件，而不能只对应一般问答正确率。

任务分解提供了另一种扩大监督范围的机会。审核长篇邮件报告时，可以分别核查每条事实的原始出处、归纳是否遗漏关键条件、行动建议是否超出权限。分解有用的条件是局部检查能够约束整体错误。如果几个局部结论分别正确，但组合顺序错误，或关键前提根本没有被列成子任务，全部局部检查通过仍不足以证明整体有效。因此，监督设计还需要检查分解是否完整，并保留最终任务层面的验证。多个评判者若共享训练数据和推理缺口，意见一致也不等于独立证据。

\subsection{评估与监督效果的检验}
\label{sec:alignment-oversight-validation}

评估体系需要分别检验行为判定是否准确，以及监督在运行中是否来得及发挥作用。规则检查器应使用已知状态变化的案例核对结果；偏好评价应保留未参与拟合的比较对并报告分歧；过程评价应与独立任务结果对照。判定器在旧数据上准确，不能自动说明新策略产生的行为也被覆盖。

监控器应评价漏报、误报、发现延迟和复核成本。其分数若用于阈值决策，还需要检查与实际偏离事件的关系。训练监控器的样本、选择阈值的数据与最终测试应分开，并包含真实使用中容易混淆的正常行为。

具体检验可以固定告警规则$m(v(h_t))\geq\tau$，在独立标注的测试轨迹上逐条运行。对实际偏离的轨迹，记录是否告警以及告警是否早于不可逆操作；对正常轨迹，记录是否误告警。由此得到检出率、误报率和提前发现比例，同时累积每次告警的人工用时。阈值$\tau$应在验证集上按允许的复核预算选择，再冻结用于测试；不能看到测试中的漏报后反复调整阈值，却仍把原测试称为独立证据。

评价时应区分三种判定。能够在少量异常样本中发现错误，支持识别能力；在真实发生率下仍能让复核队列及时处理，支持运行可用性；在损害发生前给出告警，才支持后续干预有机会生效。若两个监控器的观察范围不同，应在相同误报或复核预算下比较检出效果；否则“检出更多”可能只是因为它把几乎所有行为都送去复核。

\begin{example}[低发生率下的告警负担]
\label{ex:alignment-monitor-base-rate}
假设偏离发生率为0.1\%，监控器检出率为95\%，正常行为误报率为1\%。每10000次行为期望产生9.5次真告警、99.9次误告警，真告警仅占约8.7\%。高检出率仍可能带来大量复核。
\end{example}

监督者也需要接受故障和攻击测试。模型可能生成让评判者误判的格式，第三方内容可能污染监控上下文。可信监控不是给另一个模型加上“审核员”身份就能实现的性质，而需要独立证据、权限隔离和对失效方式的检查。

解释质量与监控效果在这里相互关联，但不能合并为一个指标。模型给出的理由即使省略了重要内部因素，也可能包含预测越权的有效信号；相反，完整描述某次局部计算，也未必能预测几步后发生的信息披露。实际测试应固定允许查看的信息，比较最终回答、动作记录和推理文本在相同误报水平下的检出能力，并检查增加信息所带来的隐私与计算成本。

还应区分旁观监控与被优化的监控。当生成模型不知道监控分数时，测试主要衡量监控器对当前行为的识别能力；当该分数被作为训练奖励或重试依据时，模型输出分布会朝更容易通过监控的区域移动。原来的漏检率不再自动成立。因此，不能只用冻结的正常样本证明一个持续反馈系统可靠。后文第~\ref{sec:security-risk-assessment}节将从主动攻击者利用反馈调整策略的角度，进一步讨论这种评估条件的变化。

\section{对齐行为的泛化与保持}
\label{sec:alignment-generalization}

在已测任务上符合要求，只说明已有条件下的表现。行为泛化关注模型面对新情境时能否继续遵守目标与约束，行为保持还包括交互变长、模型更新和系统组件变化后，原有要求是否继续成立。这里没有一种方法能够无条件保证任意未知环境中的全部行为；能够建立的结论取决于环境变化范围、可观察状态和执行机制。

维持行为有四类相互补充的机制：扩大训练与反馈覆盖，减少新情境中的学习失配；把可形式化的边界落实到执行机制，使约束不完全依赖语言模式；按完整任务管理风险与状态，避免局部小风险无界累积；在持续更新时保留已验证的能力和边界。下面先解释这些机制为何有作用，再说明怎样检验其实际范围。后节按干预位置介绍具体操作，同一种方法可能既用于即时纠偏，也参与长期保持。

\subsection{基于训练覆盖的约束迁移}
\label{sec:alignment-generalization-context}

经验泛化依赖训练中学到的是任务规则还是偶然表面模式。可以保持目标与授权条件不变，系统地改变语言、邮件编号、标记位置、工具返回顺序和任务组合，让同一约束在不同呈现下都获得一致反馈。也需要配对的边界案例：相似表述在授权条件不同、重要标记不同的情况下应得到不同动作。否则模型可能只学到某个词出现时一律拒绝，而没有学到规则依赖什么条件。

训练覆盖还应包含模型自己容易到达的错误状态。只模仿专家完整轨迹时，模型一旦在前几步偏离，就可能进入示范没有涉及的历史。数据集聚合（Dataset Aggregation, DAgger）通过反复运行当前策略、请专家为遇到的状态提供动作标签，并将数据聚合后重新学习，减少训练状态与实际访问状态的失配。Ross等在相应在线学习假设下分析了这一方法；它提供的是对策略诱导分布的性能结论，不是对任意新环境的安全承诺\scite{ross2011}

在邮件任务中，可以让候选策略在隔离环境中运行，收集漏读标记、工具返回缺失或计划中断后的历史，由评价者给出恢复读取、澄清或安全归档的示范，再加入后续训练。收集状态不要求在真实用户数据上先执行危险动作；仿真和受限接口本身也应检查覆盖差距。这里的作用机制是补上策略实际会遇到的状态，具体示范和偏好怎样更新参数将在第~\ref{sec:alignment-learning}节展开。

多情境训练、错误状态反馈和边界案例可以减少已知失配，但有限覆盖不能证明未知场景全部合规。若新环境改变了工具语义或授权规则，应先更新任务定义与状态表示，不能继续把原来的正确动作标签当成不变真值。

\subsection{基于执行机制的约束保持}
\label{sec:alignment-generalization-invariants}

对于能够形式化的行为边界，可以把约束放入不依赖模型措辞的执行机制。例如，无论用户用哪种语言要求整理邮件，执行器都根据可信授权记录核对目标对象。这样，对新表达的保护不必完全依赖模型学过类似句子；模型可以提出错误建议，但建议无法直接扩大允许动作集合。

要说明这种保持何时成立，可以采用一个简化状态模型。记$S_{\mathrm{safe}}$为满足规定约束的状态集合，$A(s)$为状态$s$下可执行的动作集合，$\mathcal T(s,a)$包含状态$s$执行动作$a$后所有被环境模型允许的后继状态。定义
\begin{equation}
 A_{\mathrm{safe}}(s)
 =\{a\in A(s):\mathcal T(s,a)\subseteq S_{\mathrm{safe}}\}.
 \label{eq:alignment-safe-actions}
\end{equation}
若初始状态属于$S_{\mathrm{safe}}$，每步都有可执行的安全动作，执行器只允许上述集合中的动作，且真实后继确实被$\mathcal T$覆盖，那么由逐步归纳可知实际状态一直留在$S_{\mathrm{safe}}$中。归纳的一步是：当前状态安全，执行器允许的动作又只能到达安全后继，因此下一状态仍安全。

\term{安全屏蔽}（shielding）在学习者与环境之间加入动作检查或修正机制。Alshiekh等依据时序逻辑规范构造屏蔽器，在动作可能违反规范时进行纠正。它说明某些约束可以通过显式环境模型与执行规则保持，而不只通过奖励惩罚促使模型遵守\scite{alshiekh2018shielding}

上述状态模型是解释不变量的教学简化。对于依赖历史的要求，状态还须包含有关历史或规范检查状态；权限、计数和外部副作用遗漏时，归纳就不能代表真实系统。新环境中的工具效果必须仍被转移关系覆盖；否则应重新验证或限制动作。安全动作集合为空时，也需要事先分析可行回退，不能假定“停止”在所有系统里都安全。该机制保证的是形式化边界，既不保证任务一定完成，也不证明模型内在目标已经对齐。具体动作校验与干预位置将在第~\ref{sec:alignment-intervention-execution}节展开。

\subsection{基于风险预算的长程控制}
\label{sec:alignment-generalization-horizon}

长过程还会积累局部风险。设$E_t$表示第$t$步发生规定偏离，无论这些事件是否独立，都有
\begin{equation}
 \mathbb P\left(\bigcup_{t=1}^{T}E_t\right)
 \leq\sum_{t=1}^{T}\mathbb P(E_t).
 \label{eq:alignment-sequential-risk}
\end{equation}
若每步只有在部署策略产生的历史下才能保持风险上界，改变策略后就要重新核查。独立性并非联合界所需条件，但局部边际概率必须针对实际行为过程成立。

因此，一个短任务中很低的失败率，不足以直接支持长任务的同一容限。若事先给定任务级容限$\alpha$，可以选择非负步骤预算$\alpha_t$，要求$\sum_{t=1}^{T}\alpha_t\leq\alpha$。若在实际策略产生的各历史下，当前步骤的条件违规概率都能由有效证据控制在$\alpha_t$以内，则可推出相应边际概率界，再由式~\eqref{eq:alignment-sequential-risk}控制任务级风险。给步骤贴上预算数字本身不会建立概率保证。任务提前结束且不再产生相关后果时，可将未执行步骤的违规事件记为空事件；已经提交、尚未完成的外部操作则仍需纳入风险范围。

例如，在一个教学设定中，任务最多100步，任务容限为0.01，均匀分配得到每步预算0.0001。该数值是待达到的要求，不是从短任务测试中自动获得的结论；若继续增加步骤，就要重新分配预算或缩小后续行动范围，不能把计数器清零后继续沿用原承诺。

机制上，可以在任务状态中保存已用次数、累计开销和剩余风险预算，在关键动作前重新检查前提，并在无法满足下一步要求时转入有明确安全条件的备用方案。缩短无效循环与减少不必要动作能够减少暴露机会，但新路线也要验证；原地停止只有在停止本身不会引发所讨论损害时才可作为安全回退。

长过程检验应保留完整任务轨迹，在不同长度和工具组合下统计至少一次违规的比例、完成任务前的中止比例和总损失。若仅保留成功完成的轨迹，就会把提前失败从统计中删掉。关键节点复核可以作为干预方案，但它是否降低整体风险，需要在同样任务与预算下比较复核前后的轨迹结果。

\subsection{持续更新中的行为保持}
\label{sec:alignment-generalization-update}

若同类问题不断出现，可以将经过核验的监督结果转化为训练反馈，修正奖励或补充边界示范。但监控器的判断并非无误标签，直接回灌可能强化误报和过度拒答。应对样本进行复核，并保留未参与更新的评价任务。

行为保持可以从训练数据、更新目标和可修改状态三个位置着手。数据回放在学习新任务时保留经过核验的旧任务与边界案例，使旧约束继续产生训练信号；输出分布约束要求新模型在保留任务上接近合规参考模型；参数约束则限制被认为对旧任务重要的权重变化。这些方法分别保留经验、行为或参数，依据不同，不能只用“微调幅度较小”概括。

弹性权重巩固（elastic weight consolidation, EWC）根据旧任务估计参数重要性，在新任务训练中更强地惩罚重要参数偏离旧值。Kirkpatrick等用这一机制研究连续学习中的遗忘。它为更新时保留既有能力提供方法，但旧任务的重要参数不等于全部安全约束，能力保持的实验结果也不能直接解释为行为对齐得到保证\scite{kirkpatrick2017ewc}

对行为约束而言，回放集合需要同时包含正常任务、易混淆的合法请求和关键违规边界，避免只回放拒答样本。参考模型若本来存在错误，保持其输出也会保留错误；参数重要性若只在常见任务上估计，也可能遗漏稀有条件。因此，保持机制仍需与独立的行为检验和执行约束共同使用。

每次干预还可能改变用户行为和未来输入。例如，大量拒答会使用户改写请求，改变监控器最初面对的分布。因此，更新的不只是模型参数，也可能是未来面对的输入分布；原先的监控阈值不能未经验证便继续沿用。

微调、模型合并、量化和持续学习都可能改变行为。参数变化小不必然意味着关键决策不变，平均能力保持也不排除少数安全场景退化。每次更新应比较原有任务、新任务、关键约束和监控器表现，并以未参与调整的数据复核。

保持对齐还需要同时评价过度拒答。若模型通过回避全部困难任务减少违规，它可能失去原本需要的服务能力。可以把合法任务完成、错误干预和约束违反分别报告，避免一个综合分数掩盖行为变化。

监督和干预也会失配：新模型或推理格式可能使旧监控器失效。因此，更新验证需要覆盖观察信息、阈值、工具权限和重试规则。

\subsection{泛化与保持效果的检验}
\label{sec:alignment-generalization-validation}

模型在熟悉提示上表现合适，不意味着新语言、组合任务或更长行动链中仍遵守同样要求。评价应区分任务难度增加造成的能力失败，与知道目标却优化了错误代理造成的目标偏离；二者可能产生相似结果，却需要不同修正。可以对同一任务构造目标说明更明确、所需知识有证据支持的对照条件，观察偏离是否减少，再分别改变语言、角色与工具条件。这样的对照有助于定位缺口，但行为失败本身通常不能唯一识别内部目标，模型声称“知道要求”也不能代替证据。

只改变任务难度还不足以覆盖情境变化。另一种更难发现的形式是：模型在通常测试条件下遵守要求，却在特定条件出现时改变行为。Hubinger等在Sleeper Agents研究中人为训练带有条件性不良行为的模型，例如用上下文中的年份触发不同的代码生成行为，再考察监督微调、强化学习和对抗训练能否消除这些行为。若干设置下，条件性行为在安全训练后仍然存在；这是人为构造的概念验证，不能据此推断普通部署模型自然形成同类策略的发生率\scite{hubinger2024sleeper}

这种在特定条件出现时激活异常行为的机制，也可能被攻击者有意植入模型；后文第~\ref{sec:security-threats}节将从后门攻击的角度讨论它。本章关心的是由此暴露的验证缺口。若评价只覆盖没有触发条件的输入，训练前后都可能得到良好分数，进而把“未激活”误认为“已消除”。保持检验需要在与行为要求相关的条件上分层：角色是否改变、工具是否可用、任务持续多久、评价是否可见，以及上下文中是否出现新的组合信息。

测试情境应尽量保留真实任务所需的行动机会。若被测模型没有工具权限，无法出现越权调用，并不说明有权限时也会主动守住边界。反过来，为了测试某种偏离而额外告知模型“当前可以违规”，又可能人为改变了原任务。有效压力测试需要记录这些差别，将构造条件下的失败与真实分布中的频率分别报告。

构造迁移测试时，可以从同一封邮件派生成对任务：仅改变请求语言、仅改变重要标记的位置、或仅增加一个会返回额外信息的工具。保持正确动作与保留约束不变，比较原条件和新条件下的任务完成、违规及监控漏报。如果某组失败增加，再检查提供缺失事实后是否恢复、执行边界是否仍有效。这样的逐因素对照有助于定位失效条件；真实部署还要覆盖因素共同变化的组合，并按实际任务分布估计总体风险。

更新后的测试也应覆盖已有纠错是否仍然有效。以邮件助手为例，可以把过去核验过的正常任务、易混淆但合法的任务和曾造成损失的任务保留为回归集合，同时补充未见过的新表达。只反复使用固定题目会逐渐把评估变成训练的一部分，因此回归集合用于发现已知退化，独立留出与新的压力测试用于发现未知退化，两者回答不同问题。

更新验收可以把旧版和新版放在相同任务、工具返回与随机条件下成对运行，分别记录双方都成功、仅旧版成功、仅新版成功和双方都失败的任务。只有新版成功的任务反映收益，只有旧版成功的任务暴露遗忘或新失效；平均成功率相同可能同时隐藏这两类变化。对关键约束还应单独检查新增违规，并让旧监控器在新版轨迹上重新评价。若需要重选阈值，应另留最终测试数据，待整个“新版策略—新版阈值—执行规则”固定后再作验收。

泛化结论应说明证据的性质。跨情境训练通过留出环境上的结果支持经验迁移；执行不变量依靠状态、转移和权限假设给出条件保证；长程预算依赖局部风险界；更新保持通过回放或正则化减少遗忘，再由独立测试核验。某个机制没有覆盖的新工具、状态或人群，应保留为未验证范围，不能由其他条件下的平均分数补足。

\section{行为干预}
\label{sec:alignment-intervention}

行为干预是在目标、约束或监督证据的指导下，改变模型将采取的行为或限制其后果。它既包括部署前修正学习过程，也包括运行中的即时纠偏。前节解释保持机制为何可能延伸到新条件；本节进一步区分何时介入、改变哪个对象，以及改变后怎样验证。

下面先讨论干预时机与选择依据，再按作用位置介绍五类操作：训练阶段的数据与参数修改、推理前的输入与上下文调整、推理中的内部计算调整、生成后的候选筛选，以及行动执行与损害恢复。训练、上下文和内部计算干预改变候选如何产生，结果筛选改变哪些候选被交付，执行端决定哪些动作真正发生；它们可以组合，最终仍需按完整行为过程验收。

\subsection{干预时机与方案选择}
\label{sec:alignment-intervention-decision}

干预不必等到违规已经发生：部署前可修正反复出现的行为偏差，运行时可在证据不足、监控告警或动作触及边界时介入，事后则处理已经产生的状态变化。反复出现的系统性问题适合反馈到训练，一次请求中的信息缺口通常先在该请求中处理。

低风险的不确定情况可以请求澄清或补充信息；可纠正的输出可以重新生成；涉及权限边界的动作应由执行端限制；高损失且不可逆的行动可能需要人工接管或停止。

一个简化的决策可以说明阈值为何依赖后果。假设继续执行造成损害的概率为$p(h)$，损害成本为$c_{\mathrm{harm}}$，停止成本为$c_{\mathrm{stop}}$，且停止能够完全避免所讨论的损害，则当
$p(h)c_{\mathrm{harm}}>c_{\mathrm{stop}}$时，停止的期望损失较小。若干预不能完全避免损害，或不同动作具有不同代价，就需要重新计算。这一规则依赖风险估计有依据，不能直接把任意监控分数当成概率。

例如，在一个教学设定中，误删重要邮件的损失为100，暂停等待复核的成本为2，那么停止阈值为$p(h)>0.02$。风险估计为0.03时，继续的期望损失为3，暂停更合适；风险估计为0.005时，继续的期望损失为0.5，单按这项损失比较则无需暂停。若动作违反的是明确授权边界，应直接由执行规则拒绝，而不是因估计概率低于阈值就放行。概率阈值用于比较允许决策中的不确定后果，授权规则用于排除根本不允许的动作。

人工接管也需要可执行条件：人应得到必要上下文，有时间作出判断，并具有停止或纠正的权限。只增加一个确认按钮，而没有改善信息和责任边界，可能只是把同样的错误交给人重复。

\subsection{训练阶段的数据与参数干预}
\label{sec:alignment-learning}
\label{sec:alignment-learning-demonstration}
\label{sec:alignment-learning-policy}

这类干预在部署前或版本更新时修改训练数据、评价目标与策略参数，作用会持续到后续任务。若偏离来自缺少正确示范，可以补充具体行为；若来自代理评价遗漏，应先修正反馈或约束，再优化策略。只在错误目标上继续训练，会放大第一节讨论的代理偏差。

示范给出可模仿的路径，评价信号区分候选优劣，参数更新使某些路径更常出现。第~\ref{sec:alignment-learning-evaluation}节已经说明反馈的来源与粒度；这里关注怎样将反馈变成生成行为的变化。训练介入应保留原版本与独立验收集，以便区分改进和附带退化。

\paragraph{通过示范修正行为}

示范数据提供在具体历史下希望采取的行为。\term{监督微调}（supervised fine-tuning, SFT）使模型增加示范输出的概率，能够传递格式、任务程序和部分约束。但学习到的是数据覆盖下的行为模式，示范没有涉及的边界仍需检验。

若回答$y$由词元$y_1,\ldots,y_m$组成，示范微调最小化
\begin{equation}
 L_{\mathrm{SFT}}(\bm{\theta})
 =-\mathbb E_{(\bm{x},y)\sim D_{\mathrm{demo}}}
   \sum_{t=1}^{m}\log\pi_{\bm{\theta}}
          (y_t\mid\bm{x},y_{<t}),
 \label{eq:alignment-sft}
\end{equation}
其中$D_{\mathrm{demo}}$是示范分布，$\bm{x}$是任务上下文，$\pi_{\bm{\theta}}$是参数为$\bm{\theta}$的生成策略，$y_{<t}$为前面的词元。训练把示范中的前缀作为条件，部署时却接着模型自己的前缀继续生成。一个早期错误可能把后续生成带入示范未覆盖的状态，这解释了为什么学会模仿标准答案，还需要在完整交互上检验。

这类训练具体需要成对记录“模型当时看到了什么”和“希望它接下来输出什么”。在邮件任务中，输入可以包含用户要求和当前可见的邮件信息，示范则依次读取重要标记、调用归档工具、确认邮件仍可检索。若模型通过文本调用工具，调用名称与参数也属于要学习的输出；工具返回作为后续动作的上下文，不应被当成模型可以任意编造的答案。式~\eqref{eq:alignment-sft}降低示范序列的负对数概率，使策略更容易产生这条路径。

示范的产物是生成策略，而不是给任意候选打分的评价器。一条正确的$A$方案示范并没有直接告诉模型$A$比$B$好多少，也没有穷举所有不允许的$C$方案。要检查学到的是程序还是表面措辞，可以保留要求而改变邮件编号、标记位置和工具返回，再让模型自己完成整段交互，检查关键动作及最终状态。

\paragraph{通过评价反馈调整策略}

获得评价之后，还需要让生成策略改变选择。显式训练奖励模型并用它优化策略，与直接从成对偏好调整策略，是连接相对评价和行为概率的两条路线。两者共享前面的评价误差与覆盖限制；是否单独保留奖励模型，不会改变反馈本身的真伪。

\term{人类反馈强化学习}（reinforcement learning from human feedback, RLHF）的一条常见路线，是把“人能比较少量行为”转化为“奖励模型能给大量新行为评分”，再用该奖励优化策略。Christiano等的工作比较行为片段，以学习到的偏好预测作为强化学习奖励，连接了人的判断与策略优化。比较反馈可以比直接示范更容易获得，但标签分歧、遗漏因素和奖励模型外推都会进入最终行为，偏好拟合不等于完整的使用安全保证\scite{christiano2017}

实际学习形成循环：当前策略产生候选行为，人比较其中一部分，奖励模型更新后再指导策略。继续收集新策略产生的数据，可以减少只依赖初始行为样本的局限，但也增加反馈成本。

在一次策略更新中，先固定奖励模型，从当前策略采样完整回答或行动轨迹，并计算每个样本的奖励。更新再利用“该样本的得分比当前通常水平高多少”来调整它的生成概率：高于基准的样本获得增加概率的信号，低于基准的样本获得相反信号。奖励模型提供训练信号，策略参数承担行为改变；若更新的是奖励模型而策略不动，生成行为不会因此自动改善。后续重新采样可以看到这种概率变化实际带来了哪些新候选。

InstructGPT将这一路线用于指令跟随：用人工示范得到初始策略，收集回答排序训练奖励模型，再使用近端策略优化调整生成策略，同时约束它偏离参考模型的程度。论文还混入预训练数据上的语言建模目标，以缓解部分任务能力退化。其贡献是展示通用语言模型能够通过这种反馈流程更符合特定标注群体对指令回答的要求；这并不代表所有人群的价值已被统一，模型仍会出现事实错误和其他失败\scite{ouyang2022instructgpt}

在这里，强化学习所需要的背景可以压缩到一个关系：模型采样出回答，奖励模型评价回答，策略更新提高较高奖励回答出现的概率。近端策略优化通过限制相对于采样策略的概率比变化，减轻单次更新过大的问题。这种优化稳定性与行为可靠性属于不同判断：即使每一步更新都很平稳，反复优化一个有偏奖励，仍会逐渐得到不符合实际意图的行为。

若只追求奖励模型得分，策略可能走向奖励模型没有可靠数据支持的输出。可以以参考策略$\pi_{\mathrm{ref}}$为基准，对偏离施加KL散度惩罚。固定上下文$\bm{x}$，理想化目标为
\begin{equation}
 \max_\pi\
 \mathbb E_{y\sim\pi}[r(\bm{x},y)]
 -\beta D_{\mathrm{KL}}(\pi\|\pi_{\mathrm{ref}}),
 \qquad \beta>0.
 \label{eq:alignment-kl-objective}
\end{equation}
假设参考策略在考虑的输出上有正概率、归一化常数有限，最优分布满足
\[
 \pi^\star(y\mid\bm{x})
 \propto\pi_{\mathrm{ref}}(y\mid\bm{x})
              \exp(r(\bm{x},y)/\beta).
\]
这个关系可以通过对概率分布加入归一化约束并求一阶条件得到。它说明奖励差能够用相对于参考策略的对数概率差表示。

为了看清归一化项的作用，令
$Z(\bm{x})=\sum_y\pi_{\mathrm{ref}}(y\mid\bm{x})
\exp(r(\bm{x},y)/\beta)$，并记上式归一化后的分布为$q$。
由$q$的定义，
\[
 \log\frac{\pi(y\mid\bm{x})}{q(y\mid\bm{x})}
 =\log\frac{\pi(y\mid\bm{x})}{\pi_{\mathrm{ref}}(y\mid\bm{x})}
  -\frac{r(\bm{x},y)}{\beta}+\log Z(\bm{x}).
\]
对$y\sim\pi$取期望并整理，式~\eqref{eq:alignment-kl-objective}的目标变为
$\beta\log Z(\bm{x})-\beta D_{\mathrm{KL}}(\pi\|q)$。
由于KL散度非负，且在$\pi=q$时等于0，最优分布正是$q$。这也说明最优性的范围是所有满足条件的分布；有限参数模型未必能够表示该分布。

将最优关系反解为奖励，有
\begin{equation}
 r(\bm{x},y)
 =\beta\log\frac{\pi^\star(y\mid\bm{x})}
                     {\pi_{\mathrm{ref}}(y\mid\bm{x})}
   +\beta\log Z(\bm{x}).
 \label{eq:alignment-reward-policy}
\end{equation}
同一上下文的两个回答相减时，未知的$\log Z(\bm{x})$抵消，恰好留下偏好概率所需要的奖励差。直接偏好学习得以避开单独估计奖励，依靠的是这一抵消，而不是因为奖励误差已经不存在。

\term{直接偏好优化}（direct preference optimization, DPO）将这一关系代入偏好模型，直接优化策略。记
\[
 d_{\bm{\theta}}(\bm{x},y^+,y^-)
 =\log\frac{\pi_{\bm{\theta}}(y^+\mid\bm{x})}
              {\pi_{\mathrm{ref}}(y^+\mid\bm{x})}
 -\log\frac{\pi_{\bm{\theta}}(y^-\mid\bm{x})}
              {\pi_{\mathrm{ref}}(y^-\mid\bm{x})},
\]
则训练目标是偏好样本上的
$-\log\sigma(\beta d_{\bm{\theta}})$的平均。它避免显式训练并调用一个独立奖励模型，但仍依赖偏好数据和所采用的模型关系\scite{rafailov2023dpo}

若$d_{\bm{\theta}}$很小甚至为负，当前策略相对于参考策略仍不够偏向优选回答；损失会推动该差值增大。对$d$求导得到$-\beta\sigma(-\beta d)$，说明当前不符合偏好排序的样本会收到较强的更新信号。对已经明显符合偏好的样本，梯度逐渐减小。这一导数解释了训练作用，却不能保证每次参数更新都增加所有优选回答的绝对概率，因为概率归一化和共享参数会使不同样本相互影响。

若某个回答在参考策略中概率为0，上述对数比便没有定义；实际模型虽然通常给词元正概率，但极低概率区域仍可能缺少可靠比较。离线偏好数据还可能只覆盖旧策略的输出，而新策略在训练中转向其他回答。DPO消除了显式奖励模型与在线强化学习的一部分复杂性，仍需处理偏好噪声、分布覆盖和过度优化。参数$\beta$与参考策略共同控制学习尺度，不能仅凭它较大或较小就推断最终行为更加安全。

仍以$A$、$B$说明两条更新路线的共同点。假设候选空间只有这两个可行方案，参考策略各给概率$1/2$，取$\beta=1$，已拟合奖励为$r(A)=\log 3$、$r(B)=0$。式~\eqref{eq:alignment-kl-objective}的理想最优策略按$\pi_{\mathrm{ref}}\exp(r)$重新加权，得到$\pi^\star(A)=3/4$、$\pi^\star(B)=1/4$。显式奖励路线先拟合这个奖励差，再通过策略优化逼近相应分布。

DPO可以直接使用第~\ref{sec:alignment-learning-evaluation}节三比一的比较数据。由于参考概率相同，此时$d_{\bm\theta}=\log[\pi_{\bm\theta}(A)/\pi_{\bm\theta}(B)]$；在能够自由调整这两个概率的理想设定下，拟合比较概率$3/4$同样得到概率比3。这一计算说明DPO省掉的是独立奖励拟合与在线奖励调用，不是比较信息。若实际模型还有大量其他输出，或参数不能独立调整，两条训练路线未必得到相同策略；上述等价只针对给定条件下的目标关系。

\subsection{推理输入与上下文的干预}
\label{sec:alignment-intervention-context}

若当前偏离源于信息缺失或任务含糊，可以在模型作出决定前改变它获得的上下文，而无需立即更新参数。具体操作包括请求用户澄清冲突要求、检索可核查证据、刷新已过时的工具状态，以及把相关授权范围明确传给模型。干预的位置是本次推理的条件，效果主要作用于当前请求；它不自动修复模型在其他任务中的同类偏差。

在邮件例子中，重要标记缺失时应重新读取或请求确认，而不是让模型依据语气猜测。用户要求“清理所有旧邮件”与既有保留规则冲突时，可以先说明冲突，再取得足以确定允许动作的信息。补充的内容也需要来源和访问范围：未经核验的检索片段可能增加错误，最新授权应来自独立权限记录，不能由模型自己补写一句“已获授权”。

对于长交互，可以把当前目标、已经确认的约束、已执行动作和未决问题保存在可核查状态中，避免关键信息因上下文截断而丢失。每次高影响动作前从真实工具状态刷新这些条件。摘要有助于控制长度，但摘要遗漏约束会造成新的失配，因此关键权限和不可逆动作记录不应只依赖模型自由生成的摘要。

验证这类干预时，应在同一任务上比较补充信息前后的动作和结果，并检查模型是否实际使用了新增证据。仅把提示写得更长，或者回答中重复约束，不足以确认行为已经改变。若缺口仍无法补齐，应进入澄清、受限执行或人工接管，而不是继续以高置信措辞掩盖信息不足。

\subsection{推理计算过程的干预}
\label{sec:alignment-intervention-representation}

输入干预改变模型所见的条件，内部计算干预则直接修改生成过程中的表示；后面介绍的输出筛选是在候选已经生成后再决定是否交付。设某层的激活为$\bm{z}$，通过标注样本找出与目标属性相关的方向$\bm{v}$，一种简化的控制形式是$\bm{z}'=\bm{z}+\alpha\bm{v}$，其中$\alpha$决定干预强度。它改变后续计算，使特定行为更容易或更难出现，而不必重新训练全部参数。

Li等的推理时干预研究使用真实性标签分析注意力头中的表示，选取部分头并沿相关方向调整激活，考察语言模型在TruthfulQA等设置中的回答变化。工作展示了内部表示可以成为行为控制的操作位置，同时也报告了真实性与有用性之间的取舍。这里观察到的特定方向不能直接解释为一个跨任务通用的“诚实开关”\scite{li2023inferenceintervention}

一般的表示干预可以按“估计方向—选定位置和强度—独立测试”构造。一个教学实现是固定层和取激活的位置，收集经核验的合规与不合规行为，计算两组激活均值之差，在差向量非零时将其归一化为方向$\bm v$；也可以训练线性分类器，用其法向量作为候选方向。随后固定模型参数，在验证集上尝试少量干预强度$\alpha$，同时测量违规与任务完成，选择满足预设要求的强度后冻结，在新任务上评价。这里说明的是一类方向构造方法，并非前述论文的完整实现。

在邮件任务中，方向标签如果只是区分“回答含有安全字样”与“不含”，得到的很可能是措辞方向。更合适的标签应来自重要邮件是否实际保留、动作是否授权等结果，并控制任务难度和输出长度。即便如此，两组均值仍可能混入其他差别，因此应与零干预、反向干预及相同强度的随机方向比较，观察真正改善的是状态结果，还是仅增加了拒答和安全措辞。

后文第~\ref{sec:interpretability-local-mechanism}节将把类似的内部计算干预用于解释，其目标与这里不同：解释需要说明某个表示怎样参与原来的计算，干预需要确认改变表示后行为是否更符合要求。一个方向能够预测标签，只说明相关；干预后得分提高，则还需排除缩短回答、增加拒答等替代原因。有效验证应同时保留任务完成度，改变测试领域、提示方式和干预强度，检查效果是否只依赖局部样本。

\subsection{生成结果的筛选与修订}
\label{sec:alignment-intervention-selection}

当目标是改变交付给用户的内容，而候选生成本身没有外部副作用时，可以先生成再筛选：拒绝某个候选、要求修订，或从多个候选中选择。这里改变的是交付规则，不必更新策略参数，也不等同于保证工具动作获得了授权。

筛选器的输入是一组尚未执行的候选及其可见证据，输出是一个被选候选或“无候选可交付”。在邮件例子中，可以先产生$A$、$B$、$C$三份计划，按保留要求排除$C$，再在$A$与$B$之间比较整理效果；如果没有候选通过约束检查，就转入澄清或人工复核，而非强行选得分最高者。排序得分与约束标记在这里用途不同，必须先定义“可交付”，再讨论可交付集合中选哪一个。候选文本通过检查后，实际工具调用仍需经过执行端授权。

干预后的新输出仍需检查。反复采样直到监控器放行，可能只是找到监控器漏检的回答。若选择规则改变了输出分布，原有可靠性和安全评价需要随之更新。重试次数、筛选规则及失败后的处理，都属于被评价的系统。

从多个候选中选择最高分回答，是一种无需更新模型参数的干预，常称为best-of-$n$选择。它将生成能力与评价能力分开：基础模型提供候选，评价器决定实际交付什么。若候选中已有优质答案且评价器能正确排序，增加候选可以改善结果；若高分区恰好有系统性误判，增加候选也会更容易找到这种误判。因而采样预算与评价器必须一起验证，不能把基础模型的单次生成结果当成整个选择系统的表现。

对只作二元放行的过滤器，可以直接算出被放行样本的风险。设基础分布中违规比例为$\rho$，违规样本通过过滤的概率为$u$，正常样本通过的概率为$v$，且分母为正，则
\begin{equation}
 \mathbb P(\text{违规}\mid\text{通过})
 =\frac{\rho u}{\rho u+(1-\rho)v}.
 \label{eq:alignment-filter-risk}
\end{equation}
若$\rho=0.1$、$u=0.05$、$v=0.9$，这一构造例子中的通过后违规比例约为$0.61\%$。结果同时取决于基础分布、漏检和正常请求通过率；只报告$u$会遗漏过滤器拒绝正常任务的代价。

在独立同分布重采样、固定过滤器且最终成功得到通过样本的理想设置下，反复采样得到的输出仍服从上述条件分布，不能仅因重试次数增加就断言每个交付结果风险必然上升。不过，若每次生成都可能发生未经隔离的副作用，累计损失会增加；若后续候选依据拒绝理由自适应改写，则独立同分布条件已经失效，需要单独评估。这些区别决定重试策略究竟改变了哪一种风险。

\subsection{行动执行与损害恢复}
\label{sec:alignment-intervention-execution}

生成层面的修改可以降低提出不当动作的概率，但外部后果最终取决于执行器是否实施该动作。因而工具授权与状态恢复需要在执行端单独落实。

行为约束可以由允许动作集合实现。执行器检查模型提出的工具、对象和参数是否属于本次授权范围，不满足条件时拒绝或缩小行动。这类边界应独立于模型是否遵守自然语言指令，由执行器依据授权记录检查；后文第~\ref{sec:security-execution-defense}节还将检验它们能否抵御外部攻击。

执行器可以将模型提案解析成“动作、对象、参数”，再与当前授权记录比较。对重要邮件的删除请求，即使模型给出很高置信度也应拒绝；对收件人不在许可集合内的转发请求，同样先阻断；归档请求则检查目标文件夹与邮件标识后才提交。为了避免检查后对象状态已经改变，应在真正执行时重新确认关键状态。模型得到的是成功结果或明确的拒绝原因，不能自行修改授权表来完成原请求。

表示控制也不适合承担所有边界。它可能降低泄露概率，却不能像执行端权限检查那样直接规定某个收件人不可访问某份文件。对于不可逆操作，应让生成层的行为改善与动作层的授权相互补充，并分别给出证据。将它们放进同一个综合分数，会遮住其中一层失效后系统还剩下什么保护。

恢复措施则需要区分可逆与不可逆影响。撤销尚未执行的动作、恢复本地状态和收回已披露的信息具有不同可行性。干预机制应尽量在不可逆后果发生前生效，并报告未能阻止的残余损失。

验证执行干预时，应固定模型提出的请求，检查合法请求能否通过、越权对象或参数是否被阻断，以及阻断时是否已经产生外部副作用。验证恢复措施则要记录实际恢复的状态和仍然残留的损失。模型在文本中表示“已经撤回”不能替代状态检查；这些结果也不能用生成内容的无害率代替。

\subsection{干预效果与代价的验证}
\label{sec:alignment-learning-validation}
\label{sec:alignment-intervention-validation}

干预是否有效，需要比较同一目标和约束下实际行为的变化。训练损失下降、风险分数降低或拒答增加，都不能单独替代任务结果与约束证据。参考策略约束也不意味着参考策略本来就安全。

检验对象应是更新后的完整行为，而不是仅比较训练损失。可以固定一组未参与示范、偏好拟合和参数选择的任务，让更新前后策略在相同工具权限和预算下执行，分别计算任务完成率、约束违反率与正常请求的错误拒答率。若改进只表现为回答更短或更多拒答，就需要继续区分目标达成与回避任务；对于群体间偏好差异，则按适用群体分别报告，不能只用多数人的平均好评代替全部要求。

奖励评价还应与独立证据对照。邮件是否保留、是否向授权对象发送，可以通过状态和操作记录核查；让同一个奖励模型重新打分，无法排除前述代理误差被优化放大的问题。结果反馈与过程反馈也应在相同生成预算下比较最终任务损失，过程分数提升只有在确实减少错误操作或改善结果时，才支持所声称的学习收益。

可以把一次验收具体化为三组观测。偏好拟合使用未训练的比较对，记录比较对数损失与排序错误；它回答评价信号是否学好。任务评价让策略在留出邮件任务上自由执行，记录归档成功、误删和未授权转发；它回答行为是否改善。约束评价另外覆盖重要标记缺失、工具返回异常等情形，记录模型是否停止、请求澄清或继续高风险动作。三组结果不互相替代：第一组改善而第二组退化，提示评价外推或优化出了问题；第二组平均改善而第三组出现关键违规，也不能按平均收益直接扩大权限。容许错误率和最低任务完成度应在观察最终测试结果之前约定，并结合样本量与估计误差判断。

检验还应覆盖训练交互中的约束，而不能只检查最终策略。一个策略可能最终学会避开危险动作，却在训练交互中已经造成不可接受损害。可以利用仿真、离线数据、受限动作空间和安全参考策略降低探索风险，但仿真差距与离线数据覆盖仍是边界。学习方法只能在这些条件下逐步扩大行为范围。

对于参数不变的运行时干预，可以在相同任务上比较无干预、仅筛选、仅表示调整及完整组合，记录任务完成、实际违规、错误阻断、复核负担和总延迟。检查候选生成阶段是否已产生副作用，并把所有重试纳入预算；否则交付结果看似合规，过程仍可能造成损害。对于训练干预，另需在版本层面比较旧任务保持和新任务收益。

完整验证应记录告警何时出现、干预何时生效，以及不可逆操作是否已经发生。若监控器准确却因处理延迟未能及时阻止损害，问题出在响应机制；若及时阻断了正常任务，则需要检查判断阈值与授权条件。这样的分层记录把“发现偏离”和“纠正偏离”的贡献分开，同时保留组合系统的最终损失。

\section{前沿研究}
\label{sec:alignment-frontiers}

前沿研究进一步追问两个边界：模型能力和推理方式改变时，可见记录还能够提供多少监督信息；当系统主动寻找监控盲点时，监督与干预的组合是否仍然有效。

\subsection{可见信息的监督价值}

一个研究方向是直接评价\term{可监控性}（monitorability），即行为偏离能否从允许观察的信息中被发现。Guan等提出干预、过程与结果性质等评价方式，并比较思维链和动作信息的监控作用。其结果说明可见过程可以提供有用信号，但监控效果随模型、任务和计算设置变化，不能据此认定思维链完整揭示了内部原因\scite{guan2026monitorability}

第~\ref{sec:alignment-oversight}节已区分解释忠实性与监督效果。可监控性研究进一步需要比较这些效果如何随任务难度、可见记录和监控预算变化：若换一种推理格式就使检出率下降，原有证据便不能直接用于新的生成过程。

\subsection{监督与干预组合的抗规避能力}

另一个方向是把监控器纳入攻击目标。Terekhov等研究了解控制协议与监控模型的攻击者，在所测协议和基准中通过适应性输入绕过监控。这个结果提醒我们检查监督与重采样等干预的组合，而不能将监控器的可信身份当作不会失效的假设\scite{terekhov2026monitors}

这些前沿问题共同推动评价单位从单个回答转向完整的行为过程。一个模型可能在单次比较中更受偏好，却需要更多重试才能通过独立检查；一个监控器可能提高检出率，却造成关键任务延迟；一个干预规则可能减少输出违规，却把错误转移到工具调用。后续研究需要在明确任务、权限、预算和损失定义的条件下比较整个过程，同时分别保留各层的测量结果，才能知道改进究竟来自哪里。

对齐证据的更新也需要可追溯的行为变化描述。只写“新版本更加安全”不足以支持扩大使用范围；更有用的结论应说明哪类任务的失败减少、正常任务是否受损、测试覆盖了哪些新能力，以及哪些场景仍需要人工或执行边界。新方法可以增加证据覆盖，却不能消除前面各章所强调的适用条件。

\section{本章小结}
\label{sec:alignment-summary}

邮件助手的目标是帮助用户整理并保留需要的信息，行为约束则规定它能操作哪些对象、何时必须核查，以及哪些后果不能用效率补偿。奖励、偏好和自然语言承诺都只是实现与观察这些要求的一部分；对齐证据必须回到真实任务、实际操作和受影响对象。

评估分别观察目标达成、约束遵守、证据使用和不确定情况下的处置。规则与状态核验、人工比较、原则指导的模型评价以及过程反馈提供不同证据；运行监测将可见记录转成告警，可扩展监督尝试弥补评价者的能力限制。评价准确、告警及时与干预有效是不同要求，必须分别验证，不能由一个好评分数替代。

新环境中的保持依赖机制和条件。训练覆盖与策略访问状态的反馈减少经验失配；执行端不变量在状态模型和权限条件成立时保持形式化边界；长程预算控制整个任务的风险范围；更新时的回放与行为、参数约束减少既有能力退化。有限测试不能保证任意未知情境，这些机制也不能自动证明模型内部目标已经符合全部要求。

干预按作用位置改变数据与参数、输入上下文、内部计算、交付结果或外部动作。选择时应区分长期学习缺口与即时信息不足，并优先在不可逆后果发生前介入。干预后，原来的输入分布、监控阈值和重试过程可能发生变化，因此需要重新核对完整行为系统的收益、违规、误拒与残余损失。前面讨论的可靠性、鲁棒性与公平性为这些判断提供了不同侧面的要求。监督可以发现某项行为不符合要求，但未必说明模型为什么作出它；下一章将讨论解释行为的证据，区分可以观察到的关联与经过检验的内部作用。隐私和对抗安全两章则将继续限定信息使用与外部操纵的边界。
